\documentclass[12pt,a4paper,oneside]{article}
\usepackage[brazil]{babel}
\usepackage[utf8]{inputenc}
\usepackage{olymp}

\setcounter{problem}{4}

\begin{document}

\begin{problem}{Fibonacci}{fibonacci}{2}
 
A sequência de \textbf{Fibonacci} é uma sequência de números naturais cujos dois primeiros elementos são 0 e 1, e todos os elementos subsequentes são calculados pela soma dos dois anteriores.

Matematicamente, representando por $F(n)$ o $n$-ésimo termo da série, começando de 0, temos que:

\begin{itemize}
\vspace{-2pt}
    \item $F(0) = 0$
\vspace{-2pt}
    \item $F(1) = 1$
\vspace{-2pt}
    \item $F(n) = F(n-1) + F(n-2)$, para todo $N \geq 2$.
\end{itemize}

Assim, os 10 primeiros termos da sequência de Fibonacci são os apresentados a seguir:

\begin{center}
\begin{tabular}{l|cccccccccc}
    $n$    & 0 & 1 & 2 & 3 & 4 & 5 & 6 & 7 & 8 & 9 \\
    \hline
    $F(n)$ & 0 & 1 & 1 & 2 & 3 & 5 & 8 & 13 & 21 & 34 
\end{tabular}
\end{center}

Faça um programa para, dado um número natural $n$, calcular $F(n)$.

%\Notes
%
%Observações, se tiver.

\InputFile

A entrada contém um único número natural $N$. Restrições: $0 \leq N \leq 90$.


\OutputFile

Seu programa deve gerar apenas uma linha de saída, contendo o valor de $F(N)$.

\Notes
Preste atenção aos tipos das variáveis, pois $F(90) \approx 2^{61}$.

\Example

\begin{example}
\exmp{
8
}{
21
}%
\end{example}

\begin{example}
\exmp{
15
}{
610
}%
\end{example}

\begin{example}
\exmp{
80
}{
23416728348467685
}%
\end{example}


%Comentários sobre o exemplo, se necessário.

\end{problem}

\end{document}
